\section{Arithmetic information and resolved conditional laws}
\label{sec:conditional-laws}
There are two different conditioning operations. An arithmetic event fixes
information about prime signs before the Poisson comparison. Conditioning
on a count or another observable then divides by the probability of that
observable. Keeping the two costs separate avoids treating a rare state as
though an absolute error resolved it automatically.


The first table records the sufficient budgets obtained directly by the
scalar and dependency-graph methods; it is not a list of the strongest
available hard-field conclusions. Here $I=-\log\PP(C)$,
$c>0$ is fixed, $Y^\sharp=\lfloor e^V\rfloor$ and
$\widehat Y=\lfloor e^{\widehat V}\rfloor$; $\nu$ and $\widehat\nu$ belong
to those cutoffs. The following rows concern the dyadic domain
$I_N=[N,2N)\cap\Z$ and the scalar or dependency-graph comparisons of
Sections~4--6. Their moving-field intensities $\Lambda$ are at least one.
\begin{center}\small
\begin{tabularx}{\textwidth}{@{}L{.25\textwidth}Xl@{}}\toprule
Observable & Sufficient budget & Prime field\\\midrule
Scalar, hard & $I+\log^+\lambda\le V-c\nu$ & $\mathcal F_{Y^\sharp}$\\
Scalar, soft & $2I+\log^+\lambda\le\widehat V-c\widehat\nu$ & $\mathcal F_{\widehat Y}$\\
All exact-level counts & $I+\log\Lambda\le V-c\nu$ & $\mathcal F_{Y^\sharp}$\\
All sites, graph / hard & $I+\log\Lambda+\log(1+\Lambda)\le V-c\nu$ & $\mathcal F_{Y^\sharp}$\\
All sites, graph / enlarged & $I+\log\Lambda\le\widehat V/2-c\widehat\nu$ & $\mathcal F_{\widehat Y}$\\
All sites, graph / adapted & $\log\Lambda\le\ell_{\rm lab,N}(I)-c\nu_N(I)$ & $\mathcal F_{Y_N(I)}$\\\bottomrule
\end{tabularx}
\end{center}
The last row uses \eqref{eq:information-saddle}--\eqref{eq:information-budget}, with $0\le I\le V_N$. Exact marks and signs can be retained in all site-labelled rows. The inequalities
have fixed negative margins; equality at the boundary and optimality for
the model are not asserted. A further conditioning on an observed count
has the separate probability cost analysed below.

\paragraph{The one-factor hard-field conclusion.}
The selected-pivot theorem improves the available site-labelled range
both on prefixes and, by restriction, on dyadic blocks:
\begin{center}\small
\begin{tabularx}{\textwidth}{@{}L{.29\textwidth}Xl@{}}\toprule
Observable and domain & Sufficient budget & Prime field\\\midrule
All prefix sites, exact excesses and signs &
$I_M+\log\Lambda\le V_M-c\nu_M$ & $\mathcal F_{Y_M^\sharp}$\\
All dyadic sites, exact excesses and signs &
$I_N^{\rm info}+\log\Lambda_N\le V_N-c\nu_N$ & $\mathcal F_{Y_N^\sharp}$\\\bottomrule
\end{tabularx}
\end{center}
The first row is \cref{thm:micro-run-tv}, with
$\Lambda=(M-L)2^{-L}\to\infty$ and
$L=\lfloor\log_2M\rfloor-a_M$, $0\le a_M=o(\log M)$.
The second is \cref{cor:dyadic-micro}, with
$\Lambda_N=N2^{-L_N}\to\infty$ and
$L_N=\lfloor\log_2N\rfloor-d_N$, $0\le d_N=o(\log N)$.
Both use the exact means $2^{-L-e-2}$ on their respective integer
index sets. The dyadic conclusion is obtained by restricting a prefix
of size $4N$ without changing $L_N$ or the integer start labels.
Thus, in this growing-intensity corridor, keeping the sites does not
require the extra intensity factor of the direct hard graph bound.
This does not enlarge the prime field. The direct routes remain useful
for their critical or relative rates and for events measurable only at
a larger or adapted cutoff; no dictionary extension is implicit.

\subsection{Stable kernels and information cost}
\begin{lemma}[Stable product lift]\label{lem:stable-lift}
Let $W$ take values in a standard Borel space, let $\mathcal F$ be a
sigma-field, and let $\nu$ be a deterministic target law. If
\[
 \EE\,\dTV(\Law(W\mid\mathcal F),\nu)\le\varepsilon,
\]
then for every $\mathcal F$-measurable random element $V$,
\[
 \dTV(\Law(V,W),\Law(V)\otimes\nu)\le\varepsilon.
\]
For $C\in\mathcal F$ with $\PP(C)>0$, the conditional bound is at most
$\varepsilon/\PP(C)$, both for $W$ and for the joint product lift with $V$.
\end{lemma}
\begin{proof}
Disintegrate the actual and product laws over $\mathcal F$. For a test set,
the difference of its conditional sections is bounded by the conditional
total-variation distance. Integrate, or integrate over $C$ and divide by
$\PP(C)$. Countable configuration spaces used here have measurable kernels;
the standard-Borel formulation permits the additional recorded variables.
\end{proof}

For $C_N\in\mathcal F_{Y_N^\sharp}$ write
$I_N^{\rm info}=-\log\PP(C_N)$. The hard scalar estimate becomes
\begin{equation}\label{eq:hard-conditioned}
 \dTV(\Law(Z_{N,L}\mid C_N),\Pois(\lambda_N))
 \ll \min\{1,e^{I_N^{\rm info}}\lambda_N
           (e^{-V_N+o(\nu_N)}+N^{-1/3+\varepsilon})\}.
\end{equation}
For $C_N\in\mathcal F_{\widehat Y_N}$ the soft proof gives instead
\begin{equation}\label{eq:soft-conditioned}
 \dTV(\Law(Z_{N,L}\mid C_N),\Pois(\lambda_N))
 \ll\min\{1,e^{I_N^{\rm info}-(\widehat V_N-\log^+\lambda_N)/2
                         +o(\widehat\nu_N)}
              +e^{I_N^{\rm info}}N^{-1/3+\varepsilon+o(1)}\}.
\end{equation}
Thus the sufficient scalar budgets are
\[
 I_N^{\rm info}+\log^+\lambda_N\le V_N-c\nu_N,
 \qquad
 2I_N^{\rm info}+\log^+\lambda_N\le\widehat V_N-c\widehat\nu_N.
\]
The hard route can be preferable at a larger information cost, but it
requires measurability with respect to its smaller prime field. These
budgets measure probabilities of events, not the number of primes exposed.
For instance, specifying $k$ independent prime signs costs $k\log2$;
a rank-$k$ compatible affine cylinder has the same cost.

At small intensity the hard estimate is relative: under
$I_N^{\rm info}\le V_N-c\nu_N$ and $\lambda_N\to0$,
\begin{equation}\label{eq:conditional-one-hit}
 \PP(Z_{N,L}>0\mid C_N)=\lambda_N(1+o(1)),\qquad
 \PP(Z_{N,L}=1\mid C_N)=\lambda_N(1+o(1)).
\end{equation}
The polynomial terms are negligible since $I_N^{\rm info}=o(\log N)$.
At large intensity, either valid scalar budget gives a conditional Gaussian
limit after centering by $\lambda_N$ and dividing by $\sqrt{\lambda_N}$.
The target Poisson normal-approximation error must still be added to the
arithmetic error; the quantitative and relative forms are in
\cref{supp:sec:conditional-details}.

\subsection{Choosing the labelled cutoff after the information cost}
\label{sec:information-cutoff}
In the dependency-graph proof, retaining all sites makes two errors
compete. Increasing the prime cutoff
removes shared-prime edges, but deletes more supports. After conditioning
on an event of cost $I$, the leading admissible logarithmic intensity is
\[
 g_{N,I}(w)=\min\{D(H_N/w)-I,\ (w-I)/2\}.
\]
The information cost enters the two constraints with different slopes.
For $0\le I\le V_N$, define
\begin{equation}\label{eq:information-saddle}
 2D(H_N/w_N(I))=w_N(I)+I,\qquad V_N\le w_N(I)\le\widehat V_N,
\end{equation}
and put
\begin{equation}\label{eq:information-budget}
 \nu_N(I)=\frac{H_N}{w_N(I)},\qquad
 Y_N(I)=\lfloor e^{w_N(I)}\rfloor,\qquad
 \ell_{\rm lab,N}(I)=\frac{w_N(I)-I}{2}.
\end{equation}
The root is unique: $2D(H_N/w)-w-I$ is strictly decreasing, with
endpoint values $V_N-I$ and $-I$. It is the crossing of a decreasing
and an increasing constraint, so $\ell_{\rm lab,N}(I)$ is the maximum
of $g_{N,I}$ on this cutoff interval.

\Needspace{15\baselineskip}
\begin{proposition}[Information-adapted labelled comparison]
\label{prop:information-labelled}
Let $C_N$ be events of positive probability, put
$I_N^{\rm info}=-\log\PP(C_N)$, and assume $0\le I_N^{\rm info}\le V_N$. For integer
$L_N\ge1$ put $\Lambda_N=N2^{-L_N}\ge1$. Suppose, for fixed $c>0$,
\begin{equation}\label{eq:information-labelled-condition}
 C_N\in\mathcal F_{Y_N(I_N^{\rm info})},\qquad
 \log\Lambda_N\le\ell_{\rm lab,N}(I_N^{\rm info})-c\nu_N(I_N^{\rm info}).
\end{equation}
For every $0<c'<c$ and $0<\varepsilon<1/3$,
\begin{equation}\label{eq:information-labelled-rate}
 \dTV\left(\Law((K^s_{x,e})_{x\in I_N,e\ge0,s\in\{\pm1\}}\mid C_N),
       \bigotimes_{x,e,s}\Pois(2^{-L_N-e-2})\right)
 \ll_{c,c',\varepsilon} e^{-c'\nu_N(I_N^{\rm info})}+N^{-1/3+\varepsilon}.
\end{equation}
All site labels, signs and exact excesses are retained. The estimate
also holds after any deterministic restriction or measurable contraction
of this field.
\end{proposition}
\begin{proof}
Write $w=w_N(I_N^{\rm info})$, $\nu=H_N/w$ and $\ell=\log\Lambda_N$.
The finite signed comparison at this cutoff has leading errors
$\exp(I_N^{\rm info}+\ell-D(H_N/w)+o(\nu))$ and
$\exp(I_N^{\rm info}+2\ell-w+o(\nu))$.
By \eqref{eq:information-saddle} and
\eqref{eq:information-labelled-condition}, their exponents are at most
$-c\nu+o(\nu)$ and $-2c\nu+o(\nu)$.
Take $E=\lceil(I_N^{\rm info}+\ell+w)/\log2\rceil$. The enlarged support is
$\log_2N+O(V_N)$, uniformly in the parameters; the real conditional and
target tails are $O(e^{-w})$. The polynomial remainder is at most
$e^{I_N^{\rm info}}\Lambda_N(1+\Lambda_N)N^{-1/3+\varepsilon/2+o(1)}$.
Since $I_N^{\rm info}+2\ell\le w=O(V_N)=o(\log N)$, it is absorbed by the
stated remainder. The free-cutoff estimates are uniform on
$[V_N,\widehat V_N]$; the detailed calculation, including the support
and tail checks, is in \cref{supp:sec:information-cutoff}.
\end{proof}
At zero information cost, $w_N(0)=\widehat V_N$ and
$\ell_{\rm lab,N}(0)=\widehat V_N/2$, recovering the enlarged-cutoff
labelled range. For $0<I<V_N$, the allowed leading value is strictly
larger than $\widehat V_N/2-I$. This is a comparison of proof budgets
for events measurable at the required cutoff, not a transfer of events
from a larger sigma-field into a smaller one. If $I_N^{\rm info}/V_N\to\theta\in[0,1]$,
then
\begin{equation}\label{eq:information-leading}
 \frac{\ell_{\rm lab,N}(I_N^{\rm info})}{V_N}
       \longrightarrow\frac{\sqrt{\theta^2+8}-3\theta}{4}.
\end{equation}
The exact saddle, rather than this first-order limit, is used in
\eqref{eq:information-labelled-condition} to retain its $\nu_N(I_N^{\rm info})$
margin.

\paragraph{Admissible information matters.}
For a fixed event known only to belong to
$\mathcal F_{\lfloor e^{w_0}\rfloor}$, with
$w_0\in[V_N,\widehat V_N]$, optimize over $w\ge w_0$, not over all
cutoffs. The best leading budget on this interval is attained at
$w=\max\{w_0,w_N(I)\}$. The sufficient condition there is
$\log\Lambda_N\le g_{N,I}(w)-cH_N/w$ with that event measurable in
$\mathcal F_{\lfloor e^w\rfloor}$. The proof is the same finite
comparison. These optimized upper bounds do not identify an intrinsic
phase boundary of the multiplicative model, nor do they use the scalar
soft-exception estimate as a field theorem.

The exact prefix comparison of \cref{thm:micro-run-tv}, and its dyadic
restriction \cref{cor:dyadic-micro}, use the hard field. Deep starts and
signed tails are checked independently in
\cref{supp:pivot:lem:inputs}. Its partition readouts require no mesh
hypothesis. Process limits still require their own target, ordering and
censoring estimates. Neither that contraction
nor this information-adapted graph calculation silently changes the
admissible prime field of the other route.

\subsection{The exact resolution cost}
\begin{lemma}[Sharp conditioning inequality]\label{lem:conditioning-floor}
Let $\mu,\nu$ be probability laws on $E\times F$ with
$\dTV(\mu,\nu)\le\varepsilon$. Let $B\subset E$ have
$p=\nu_E(B)>\varepsilon$ and write $q=\mu_E(B)$. Then $q>0$ and
\begin{equation}\label{eq:sharp-conditioning}
 \dTV(\mu_F(\cdot\mid B),\nu_F(\cdot\mid B))
 \le\frac{\varepsilon}{\max(p,q)}\le\frac{\varepsilon}{p}.
\end{equation}
\end{lemma}
\begin{proof}
The bound $|p-q|\le\varepsilon$ gives positivity. First suppose $q\le p$.
For any measurable $D\subset F$ put $A=B\times D$. Since $\mu(A)\le q$,
\[
 \frac{\mu(A)}q-\frac{\nu(A)}p
 \le\frac{\mu(A)-\nu(A)+p-q}{p}
 =\frac{\nu(B\times D^c)-\mu(B\times D^c)}p
 \le\frac\varepsilon p.
\]
The supremum of the positive differences equals total variation for
probability laws. If $q\ge p$, interchange the two laws.
\end{proof}
The denominator expresses the correct scale: convergence after conditioning
requires $\varepsilon/p\to0$, not merely $p>\varepsilon$.
This lemma applies to an entire future path at once.

Let $\lambda_{N,m}=2^{\vartheta_N-m}$ and
$p_{N,m}(n)=e^{-\lambda_{N,m}}\lambda_{N,m}^n/n!$.
Write $\varepsilon_{N,m,C}^{\rm path}$ for the total-variation error between
the actual complete future threshold path, under $C_N$, and the target
staircase based at $m$. The aggregated proof works equally for an arbitrary
base threshold in a fixed logarithmic band (including $\lambda_{N,m}<1$).
In its nontrivial information range it gives
\begin{equation}\label{eq:path-error}
 \varepsilon_{N,m,C}^{\rm path}
 \ll e^{I_N^{\rm info}}\lambda_{N,m}
       (1+\log^+(2\lambda_{N,m}))e^{-V_N+o(\nu_N)}
       +N^{-1/3+\varepsilon},
\end{equation}
with an arbitrarily small fixed exponent slack absorbing the enlarged
length and the subpolynomial information factor.

\begin{theorem}[Resolved thinning kernels]\label{thm:resolved-kernel}
If $p_{N,m}(n)>\varepsilon_{N,m,C}^{\rm path}$, then
\begin{equation}\label{eq:resolved-kernel}
 \dTV\left(\Law((Z_{N,b_N+m+j})_{j\ge0}\mid C_N,Z_{N,b_N+m}=n),
                     \operatorname{Thin}_n\right)
 \le\frac{\varepsilon_{N,m,C}^{\rm path}}{p_{N,m}(n)},
\end{equation}
where $\operatorname{Thin}_n$ is the path of $n$ independent geometric
lifetimes from \eqref{eq:target-forward}. If a fixed lower segment is also
retained, the same inequality uses the comparison error based at its
lowest level and gives the reverse immigration kernel
\eqref{eq:target-reverse} jointly with the future.
\end{theorem}
\begin{proof}
Apply \cref{lem:conditioning-floor} to the present count and the entire
retained path. The conditional target is exactly the geometric-lifetime
path by independent Poisson splitting. For lower levels, the increments
below the present level are independent Poisson variables. A finite
intensity half-line comparison cannot control the whole infinite reverse
path, whose total target mass diverges.
\end{proof}
For every integer $n\ge1$ and every $\lambda>0$, the local formula is
\[
 p_\lambda(n)=\frac{e^{-\delta_n}}{\sqrt{2\pi n}}
                 e^{-\lambda h(n/\lambda)},\qquad
 0\le\delta_n\le\frac1{12n},\qquad h(v)=v\log v-v+1.
\]
The correction is independent of $\lambda$; its relative error is at most
$1/(12n)$. For $n=0$ use the exact mass $e^{-\lambda}$ instead.
The telescoping proof is in \cref{supp:sec:conditional-details}.
For $n_N\to\infty$, a sufficient condition for the right side of
\eqref{eq:resolved-kernel} to vanish is
\begin{equation}\label{eq:resolved-budget}
 I_N^{\rm info}+\log^+\lambda_{N,m}
 +\log(1+\log^+(2\lambda_{N,m}))
 +\tfrac12\log n_N+\lambda_{N,m}h(n_N/\lambda_{N,m})
 \le V_N-c\nu_N.
\end{equation}
In the central regime $n_N=\lambda_{N,m}+t_N\sqrt{\lambda_{N,m}}+O(1)$,
$|t_N|=o(\lambda_{N,m}^{1/6})$, this reduces, with a fixed margin, to
\[
 I_N^{\rm info}+\tfrac32\log\lambda_{N,m}+\tfrac12t_N^2
 \le V_N-c\nu_N.
\]
The coefficient $3/2$ is a sufficient error budget, not a threshold of the
model. Its derivation, including the common-band check, is in
\cref{supp:sec:conditional-details}.

\subsection{Quenched comparison on geometric scales}
\begin{corollary}[Almost-sure environmental control]\label{cor:quenched}
Let $N_k\asymp q^k$, $q>1$, and suppose
$\log\Lambda_{N_k}\le V_{N_k}-c\nu_{N_k}$.
For each $0<\beta<c$, the conditional total-variation distance, given
$\mathcal F_{Y_{N_k}^\sharp}$, between the complete aggregated staircase
and its deterministic target is at most $e^{-\beta\nu_{N_k}}$ for all
sufficiently large $k$, almost surely.
\end{corollary}
\begin{proof}
The fibrewise proof of \cref{thm:moving-comparison}(ii) gives mean distance
$e^{-c\nu_N+o(\nu_N)}+N^{-1/3+o(1)}$.
Markov's inequality at $e^{-\beta\nu_N}$ gives a summable series along
$N_k\asymp q^k$, since $\nu_{N_k}\asymp\sqrt{k/\log k}$.
Borel--Cantelli concludes; no independence between scales is used.
\end{proof}
The same argument applies to the dyadic graph-based labelled field in
the sufficient range of that estimate, and to scalar kernels uniformly over the $O(\log N)$
lengths of any fixed logarithmic band with a fixed error margin.
A maximal coupling exists separately in each qualified environment and
scale. This does not furnish a canonical coupling over all scales.
